% Chapter 5 -- Rubric: PO4 (4.1)
\chapter{Investigation and Evaluation}
\label{ch:evaluation}

\section{Experimental Investigation}
\label{sec:investigation}

\subsection{What had to be established, and in what order}

The obvious question to ask of an exploratory testing agent is how many defects
it finds. This project cannot answer that question, for a reason set out in
\Cref{subsec:measurement} and repeated here because it governs the whole
chapter: computing precision or recall requires a build whose defects are known
in advance, and the partner supplied neither a defect export nor a seeded-defect
build (constraint C3). No system in the reviewed literature reports those
figures either, so the gap is the field's and not only this project's.

What can be established is the question that comes before it, and that the
literature largely skips. A failed test run has three possible explanations:
the application is wrong, the agent is wrong, or the environment broke. Until
those are separated, a defect count is uninterpretable, because ``the test
failed'' and ``the agent got lost'' are the same observation. The investigation
was therefore ordered around three questions, each of which had to be answered
before the next was worth asking.

\begin{description}[leftmargin=2.4em,style=nextline]
  \item[RQ1. Can the agent complete a campaign without human rescue?] The
    measure is \emph{autonomy}: the proportion of executions whose outcome is a
    statement about the application rather than a failure of the agent or its
    environment. This is the precondition for every other measurement.
  \item[RQ2. Does accumulated knowledge make later tests cheaper and better
    directed?] The system's central design claim is that a persistent graph
    shared between agents improves the next test. The measures are median
    device steps per test and the growth of mapped interface states.
  \item[RQ3. Does the system transfer to applications and platforms it was not
    built against?] The measure is whether a new target requires code, and what
    the results look like when it does not.
\end{description}

\subsection{Method}

Campaigns were run as a controlled sequence. Each campaign changed one property
of the agent and was then measured against the same application and the same
specification, so that a movement in the result could be attributed to the
change rather than to the target. Ninety-eight device executions were recorded
this way, each preserved on disk as a complete step-by-step trajectory, which
is what makes the sequence auditable after the fact rather than only at the
time (NFR-D5).

Two properties of the design make the measurement possible at all. The first is
the fault-attribution taxonomy of \Cref{subsec:attribution}: every execution is
classified at the point of failure as application evidence, agent limitation or
environment fault, by the system rather than by a reader afterwards. The second
is the degradation register, which records every fallback and every silent
substitution, so that a campaign that appears to have succeeded because a
component quietly degraded can be identified as such (NFR-D2).

The honest limitation of the method is stated here rather than buried. The
model assignment did not stay fixed across the sequence (\Cref{tab:modelhistory}),
so where a campaign changed both an agent mechanism and a model, the attribution
of the improvement between them is not clean. This is noted against the rows
where it applies.

\subsection{RQ1: autonomy, and why it is the headline}

\Cref{tab:campaigns} is the campaign sequence against the primary Android
target. It is the central experimental result of the project.

\begingroup
\small
\setlength{\tabcolsep}{5pt}
\begin{xltabular}{\textwidth}{@{}>{\hsize=0.52\hsize}Y>{\hsize=0.20\hsize}Y>{\hsize=0.26\hsize}Y>{\hsize=1.04\hsize}Y>{\hsize=0.98\hsize}Y@{}}
  \caption[Campaign sequence against the primary target]{The campaign sequence against the primary Android target. One property of the agent changed per campaign; the application and the specification were held constant throughout. Reconstructed from the on-disk trajectory directories, one per device execution.}
  \label{tab:campaigns}\\
  \toprule
  Campaign & Tests & Wall clock & What changed & Outcome \\
  \midrule
  \endfirsthead
  \caption[]{Campaign sequence \emph{(continued)}}\\
  \toprule
  Campaign & Tests & Wall clock & What changed & Outcome \\
  \midrule
  \endhead
  \midrule
  \multicolumn{5}{r@{}}{\footnotesize\itshape continued on the next page}\\
  \endfoot
  \bottomrule
  \endlastfoot

  \mbox{24 Jul, 22:25} & 6 & 67 min & First end-to-end wiring & Smoke run, taken before instrumentation existed. No attribution available. \\
  \addlinespace
  \mbox{25 Jul, 01:42} & 4 & 14 min & Graph write-back & Smoke run. Confirmed that an execution reaches the graph. \\
  \addlinespace
  \mbox{12 Aug, 22:42} & 25 & 107 min & First full campaign with attribution & \textbf{0\% pass.} 82\% of device steps were unparseable, so almost nothing could be attributed. The instrumentation worked; it revealed that the evidence did not. \\
  \addlinespace
  \mbox{13 Aug, 03:03} & 29 & 84 min & Non-reasoning vision model; livelock detector & \textbf{20\% pass, autonomy 96\%.} Two changes at once, so the split between them is not established. \\
  \addlinespace
  \mbox{13 Aug, 19:02} & 9 & 16 min & (none) & \textbf{Cancelled deliberately} after bad preconditions and unreadable logs were spotted in progress. \\
  \addlinespace
  \mbox{13 Aug, 19:55} & 25 & 61 min & Precondition filter, containment matching, device reset, log reassembly & \textbf{36\% pass, 100\% autonomy, zero silent fallbacks.} Not one execution was lost to a failure of the agent. \\
  \bottomrule
\end{xltabular}
\endgroup

The result that matters in that table is not the pass rate. It is the autonomy
column: a figure that was effectively zero in the 12 August campaign reached
100\% on 13 August, meaning every one of twenty-five executions ended in a
statement about the application rather than in the agent getting lost. Before
that point no defect number the system produced could have meant anything. The
sequence of fixes that achieved it is worth naming because none of them is
glamorous: log reassembly so that device steps could be parsed at all,
containment matching so that a screen could be recognised when its identifiers
were absent, a precondition filter so that tests were not issued against states
the application could not be in, and a device reset so that one test could not
contaminate the next.

The fifth campaign is included deliberately although it produced no data.
Inspecting a run in progress and stopping it once the configuration is known to
be bad is cheaper than letting it generate twenty-five misleading results, and
a table that recorded only the campaigns that finished would misrepresent how
the work was actually done.

\subsection{RQ2: whether accumulated knowledge pays}

The graph is the system's central bet, and it is measurable through the route
the executor takes. Across the campaign sequence the \emph{median number of
device steps per test fell from 50 to 9}. The mechanism is the navigation
memory of ETA-REQ-302: paths that previously reached a screen are recorded,
along with the transitions that led nowhere, so a later test walks a known
route instead of rediscovering it. At the end of the sequence the navigation
memory held 123 nodes and 95 recorded avoid-transitions.

This is the clearest evidence in the project that the knowledge graph does what
it was built to do, and it is also the largest single cost saving achieved,
because the device agent accounts for roughly 92\% of the cost of a test
(\Cref{sec:budget}) and that cost is very nearly linear in steps. Two other
accumulations are visible in the same runs: mapped interface states grow
monotonically within a campaign, reaching 97 distinct states on the primary
target, and mined error patterns feed a strategy memory that the planner reads
when choosing what to attempt next (ETA-REQ-303).

The claim this evidence does not support is that the graph improves the
\emph{quality} of test selection, as opposed to the efficiency of execution.
Step count and state coverage are both execution measures. Whether a test
proposed from a rich graph is a better test than one proposed from an empty
graph is exactly the question that needs ground truth to answer, and it is
carried into \Cref{sec:conclusion} with the seeded-defect experiment.

\subsection{RQ3: transfer to targets the system was not built against}

Five targets were exercised in September across both platforms.
\Cref{tab:targets} reports them together. None required a code change: each was
added as a target profile, which is requirement DR-03 and DR-09 in operation
rather than in principle.

\begingroup
\small
\setlength{\tabcolsep}{4pt}
\begin{xltabular}{\textwidth}{@{}>{\hsize=0.62\hsize}Y>{\hsize=0.34\hsize}Y>{\hsize=0.28\hsize}Y>{\hsize=0.50\hsize}Y>{\hsize=0.42\hsize}Y>{\hsize=0.84\hsize}Y@{}}
  \caption[Results across all five targets]{Results across all five targets exercised in September. ``Agent-attributed'' is the share of executions lost to the agent rather than resolved against the application; lower is better, and it is the figure that decides whether the rest of the row can be read at all.}
  \label{tab:targets}\\
  \toprule
  Target & Platform & Tests & App-level failures & Agent-attributed & Coverage \\
  \midrule
  \endfirsthead
  \caption[]{Results across all five targets \emph{(continued)}}\\
  \toprule
  Target & Platform & Tests & App-level failures & Agent-attributed & Coverage \\
  \midrule
  \endhead
  \midrule
  \multicolumn{6}{r@{}}{\footnotesize\itshape continued on the next page}\\
  \endfoot
  \bottomrule
  \endlastfoot

  DataGhurhi (departmental form and survey platform) & Web & 64 & 4 & \textbf{0\%} & 51 passed, 7 skipped, 2 stopped by a cap \\
  \addlinespace
  AUT-1, the production marketplace & Android & 45 & 11 & 31\% (14) & 33/106 requirements, 29 feature areas, 97 UI states \\
  \addlinespace
  Google Contacts 4.86.23 & Android & 50 & 19 & 20\% (10) & 55/141 requirements, 35 feature areas, 60 UI states \\
  \addlinespace
  broken-workshop & Web & 10 & 0 & 40\% (4) & 170 steps; 5 runs completed \\
  \addlinespace
  AcademyBug & Web & 20 & 3 & \textbf{80\% (16)} & 0 of 0 requirements; no specification ingested \\
  \bottomrule
\end{xltabular}
\endgroup

\paragraph{DataGhurhi, the strongest result the project produced.} Fifty runs
against a departmental form and survey platform yielded 64 judged test cases,
of which 51 passed and 4 were product failures, with \emph{nothing at all
attributed to the agent}. A campaign in which the tester never once becomes the
explanation for a failure is the condition the whole of \Cref{tab:campaigns}
was working towards, reproduced here on a platform and a domain the tuning was
not done against. The four findings were one critical and three major: a
notifications control absent from the workspace header, and the Quantitative
Analysis sidebar control failing to open its view. The report is filtered to
exclude agent-caused failures before a defect is ever named, which is the
fault-attribution taxonomy acting as a gate on what reaches a stakeholder
(DR-08).

\paragraph{Google Contacts, the strongest generalisation result.} This target
matters for a different reason: it is an application no member of the team had
any hand in, tested against an independently written 141-requirement
specification rather than a vendor document. The agent designed 50 test cases
for an application it had never seen, executed each on a device, and returned
19 application-level failures and 125 findings about the application, covering
55 of 141 requirements across 35 feature areas in 6.2 hours of device time at a
median of 346 seconds per test. Nothing about the target was known to the
system in advance except its specification, which is the claim DR-01 and DR-03
make together.

\paragraph{AUT-1, the hardest target.} The production marketplace is the target
the partner cared about and the one the system found most difficult, for
reasons that are structural rather than incidental. Its cross-platform toolkit
exposes almost no stable identifiers through the accessibility tree and reports
a single activity on every screen, so screens must be distinguished
structurally (constraint C4). It also carries live users, so registration,
one-time-password verification, payment and account deletion were blocked in
configuration rather than exercised (C5, DR-07). Within those limits the
campaign returned 11 application-level failures and 95 findings about the
application across 45 tests, mapping 97 distinct interface states, the most of
any target.

\paragraph{The two results that did not go well.} They are reported here rather
than in a footnote because a table containing only successes would not be
evidence of anything. On \emph{broken-workshop}, a deliberately faulty web
application, the agent executed ten tests, completed five, and found
\emph{zero} defects. Either the reachable faults were not the ones the tests
targeted, or the agent did not provoke them; the campaign cannot distinguish
these, and it is a direct illustration of why recall cannot be claimed without
ground truth. On \emph{AcademyBug}, 16 of 20 runs were lost to the agent's own
limits, an 80\% agent-attribution rate that is the exact inverse of the
DataGhurhi result. The cause is identified in \Cref{sec:revisions} and it is
instructive: no specification was ingested for that target, so the coverage
line reads ``0 of 0 requirements''.

\subsection{What the investigation establishes}

Autonomy is demonstrated and reproduced: from effectively zero to 100\% on the
primary Android target, and independently to 100\% on a web target in a
different domain. Knowledge accumulation is demonstrated in execution
efficiency, with median device steps falling from 50 to 9. Transfer is
demonstrated across five targets, two platforms and three domains with no code
change. Unit cost is measured rather than estimated, at approximately USD 0.008
per executed test (\Cref{sec:budget}).

What is not established is the value of the output. Every application-level
failure in \Cref{tab:targets} is a \emph{candidate} defect requiring human
confirmation, and the system is scoped to say so. Precision and recall are not
reported, and on the evidence of \Cref{subsec:measurement} they are not
reported by the comparable literature either. A reader should take the cost and
autonomy results as established and treat the defect counts as what they are:
the output of a tool whose accuracy has not yet been measured against a known
answer.

\section{Performance Evaluation}
\label{sec:evaluation}

\subsection{Setup and acceptance criteria}

All campaigns ran the four-process stack of \Cref{fig:topology} on a single
developer laptop: gateway, retrieval service, executor and a Neo4j server, with
an Android emulator or a browser alongside. Android targets were driven through
the accessibility service using mobilerun with its companion on-device
application; web targets were driven through an accessibility-annotated browser
snapshot. Models were as configured in \Cref{tab:models}. Each execution was
bounded by an interaction budget and a wall-clock limit, the values used for
these campaigns being those recorded in \Cref{tab:budgets}.

A requirement is recorded below as \emph{met} only where a campaign artefact
demonstrates it, not where the code implementing it exists. Where a requirement
is implemented but could not be exercised for want of an input the project did
not control, it is recorded as \emph{met, unexercised}, and the missing input is
named.

\subsection{The partner's requirements}

\Cref{tab:eta-eval} evaluates the eight requirement families of
ETA-TR-2026-001~\cite{etatr}. All eight have working implementations verifiable
on the dashboard, and the ten work packages derived from them were built,
integrated and verified.

\begingroup
\small
\setlength{\tabcolsep}{5pt}
\begin{xltabular}{\textwidth}{@{}>{\hsize=0.46\hsize}Y>{\hsize=0.80\hsize}Y>{\hsize=0.40\hsize}Y>{\hsize=1.34\hsize}Y@{}}
  \caption{Evaluation against the partner's eight requirement families}
  \label{tab:eta-eval}\\
  \toprule
  Requirement & Subject & Status & Evidence \\
  \midrule
  \endfirsthead
  \caption[]{Evaluation against the partner's requirements \emph{(continued)}}\\
  \toprule
  Requirement & Subject & Status & Evidence \\
  \midrule
  \endhead
  \midrule
  \multicolumn{4}{r@{}}{\footnotesize\itshape continued on the next page}\\
  \endfoot
  \bottomrule
  \endlastfoot

  ETA-REQ-301 & Defect history as a knowledge source & Met, unexercised & Defect nodes, ingestion endpoint and defect-weighted retrieval are live, and a defect-history block is injected into generation. No campaign exercised it: the partner supplied no defect export (C3). \\
  \addlinespace
  ETA-REQ-302 & Execution-path memory as a navigation tree & Met & Median device steps per test fell from 50 to 9 across the campaign sequence. Navigation memory held 123 nodes and 95 avoid-transitions. \\
  \addlinespace
  ETA-REQ-303 & Incremental experiential learning & Met & Error-pattern mining and strategy memory populated from execution logs and read by the planner; three mined patterns present at the end of the sequence. \\
  \addlinespace
  ETA-REQ-304 & Multi-dimensional partitioning by profile, platform and application & Met & Five targets across two platforms held in one graph without cross-contamination (\Cref{tab:targets}). \\
  \addlinespace
  ETA-REQ-305 & Self-healing and adaptive execution & Met & Failure classification with per-category retry; the livelock detector introduced in the fourth campaign and the device reset in the sixth are both in this family. \\
  \addlinespace
  ETA-REQ-306 & Regression risk scoring and risk-weighted prioritisation & Met & Per-area risk scores drive prioritisation and are exposed through the risk API and the dashboard. \\
  \addlinespace
  ETA-REQ-307 & Test-case quality metrics and feedback loop & Met & Per-test effectiveness recorded and fed back into generation; duplicate suppression operates against executed-test titles. \\
  \addlinespace
  ETA-REQ-308 & Anomaly detection over execution patterns & Met & Failure-rate spikes, duration regressions and new-failure detection implemented and surfaced on the dashboard. \\
  \bottomrule
\end{xltabular}
\endgroup

\subsection{The team's derived requirements}

\Cref{tab:dr-eval} evaluates the eleven derived requirements and the six
non-functional requirements of \Cref{ch:requirements}. Two are recorded with
qualifications, and those qualifications are the substance of
\Cref{sec:revisions}.

\begingroup
\small
\setlength{\tabcolsep}{5pt}
\begin{xltabular}{\textwidth}{@{}>{\hsize=0.24\hsize}Y>{\hsize=0.86\hsize}Y>{\hsize=0.42\hsize}Y>{\hsize=1.48\hsize}Y@{}}
  \caption{Evaluation against the derived and non-functional requirements}
  \label{tab:dr-eval}\\
  \toprule
  Req. & Subject & Status & Evidence \\
  \midrule
  \endfirsthead
  \caption[]{Evaluation against the derived requirements \emph{(continued)}}\\
  \toprule
  Req. & Subject & Status & Evidence \\
  \midrule
  \endhead
  \midrule
  \multicolumn{4}{r@{}}{\footnotesize\itshape continued on the next page}\\
  \endfoot
  \bottomrule
  \endlastfoot

  DR-01 & No knowledge source is a hard dependency & Met, with a caveat & AcademyBug ran with no specification at all and produced findings, so the dependency is genuinely soft. But 80\% of its runs were lost to the agent, against 0\% on DataGhurhi where a specification was present. The system runs without one; it does not run well. \\
  \addlinespace
  DR-02 & Learn the application's structure from the running application & Met & 97 distinct interface states mapped on AUT-1 and 60 on Contacts, from execution alone, with no design export. \\
  \addlinespace
  DR-03 & Application-agnostic by construction & Met & Five targets, three domains, two platforms, no code change; each added as a target profile. \\
  \addlinespace
  DR-04 & A second execution platform & Met & Three web targets executed through the browser executor against the same planner and the same graph. \\
  \addlinespace
  DR-05 & Evidence integrity & Met & The 12 August campaign recorded 82\% of device steps as unparseable; after log reassembly the sixth campaign recorded zero silent fallbacks. \\
  \addlinespace
  DR-06 & Turn trajectories into durable, atomic knowledge & Met, with a caveat & 163 and 198 findings produced on the two Android campaigns, each attributed and deduplicated. Dedup is imperfect: DataGhurhi returned three near-identical findings for one control. \\
  \addlinespace
  DR-07 & Enforced safety boundary & Met & Registration, one-time-password verification, payment and account deletion blocked at the configuration layer on the production target. No irreversible action occurred in 98 executions. \\
  \addlinespace
  DR-08 & An operator surface & Met & A campaign is configured, launched and read from the dashboard, and exports a stakeholder-facing report that filters agent-caused failures out before naming a defect. \\
  \addlinespace
  DR-09 & Configuration as data & Met & Every target in \Cref{tab:targets} was added as configuration; every model in \Cref{tab:models} is a settings value. \\
  \addlinespace
  DR-10 & Measured, not estimated, cost & Met & USD 0.008 per executed test, derived from metered billing and 232 stored observations, not from a price card (\Cref{sec:budget}). \\
  \addlinespace
  DR-11 & The planner must be able to investigate, not only receive & Met, not default & Both planner modes are implemented behind one output contract. The retrieval-issuing mode is the default; the tool-calling mode remains switched off pending a campaign that compares them. \\
  \addlinespace
  NFR-D1 & Cost ceiling per executed test & Met & USD 0.008, roughly one taka, against published figures of USD 10 to 22 for comparable systems (\Cref{sec:feasibility}). \\
  \addlinespace
  NFR-D2 & Observability of every degradation and fallback & Met & Degradation register live; the sixth campaign is recorded as having zero silent fallbacks, which is a claim the register makes checkable. \\
  \addlinespace
  NFR-D3 & Attribution soundness & Met & The central result of \Cref{sec:investigation}: 100\% autonomy on two independent targets, and every reported failure carries its attribution class. \\
  \addlinespace
  NFR-D4 & No irreversible action on the application under test & Met & Enforced in configuration, not by convention; no violation in 98 executions across five targets. \\
  \addlinespace
  NFR-D5 & Reproducibility of a campaign & Met, with a caveat & All 98 executions preserved as complete trajectories on disk. Campaign-to-campaign comparison is nevertheless limited by the model changes in \Cref{tab:modelhistory}. \\
  \addlinespace
  NFR-D6 & Confidentiality of test credentials & Met & The application password reaches the executor goal only; the planner prompt carries the role name and never the credential. \\
  \bottomrule
\end{xltabular}
\endgroup

\section{Deviations and Design Revisions}
\label{sec:revisions}

Six deviations are recorded. For each, the cause was diagnosed from the
trajectory record rather than inferred, and the revision is named. Two were not
revised, and the reason is given.

\paragraph{D1. Evidence was unreadable before it was wrong.} The first
instrumented campaign returned a 0\% pass rate, which initially read as a
catastrophic quality result. It was not: 82\% of device steps could not be
parsed, so almost no execution could be attributed at all. The cause was that
step records were being fragmented across log boundaries and reassembled
incorrectly. \emph{Revision:} log reassembly at the gateway, so a step is
written once and completely. The campaign that followed recorded zero silent
fallbacks. This deviation is the origin of NFR-D2 being treated as a
requirement rather than a convenience.

\paragraph{D2. The agent's own failures were being reported as application
defects.} Early runs counted a navigation failure and an exhausted interaction
budget as evidence about the application. Four executions had additionally been
recorded as agent navigation faults when the browser process had simply ceased
to exist. \emph{Revision:} the three-way fault-attribution taxonomy of
\Cref{subsec:attribution}, applied at the point of failure and enforced in the
report generator, so that an agent limitation cannot reach a stakeholder as a
defect. The dead-browser condition was reclassified as an environment fault
that ends the batch. This is NFR-D3, and it is the precondition for every
number in \Cref{tab:campaigns}.

\paragraph{D3. The loop stalled without terminating.} Executions consumed their
full wall-clock budget while making no progress, oscillating between two
screens. \emph{Revision:} a livelock detector, introduced in the fourth
campaign, which terminates a run that is not advancing and records it as an
agent limitation rather than letting it consume the budget silently. The
AcademyBug trajectories show the untreated form of this failure clearly: runs
cycling between \modpath{find-bugs} and \modpath{report-bugs} until they timed
out.

\paragraph{D4. Screens could not be told apart on the primary target.} The
cross-platform toolkit used by AUT-1 exposes almost no stable resource
identifiers and reports a single activity for every screen, so the original
identity strategy, which assumed identifiers, collapsed (constraint C4).
\emph{Revision:} structural state identity with containment matching, so a
screen is recognised by the shape and content of its accessibility tree rather
than by a name. This is what made the 97 mapped states on that target possible.

\paragraph{D5. Quality degrades sharply without a specification, and this was
not designed for.} DR-01 requires that no knowledge source be a hard
dependency, and the requirement is met in the literal sense: AcademyBug was run
with nothing ingested and still produced three defects. But 16 of its 20 runs
were lost to the agent, against none of DataGhurhi's 64. Without a
specification there are no requirements to target, so the planner has no
stopping condition for a test and no criterion for what the application ought
to have done; the coverage line reads ``0 of 0 requirements'' because there was
nothing to cover. \emph{Not revised.} The correct revision is a generated
provisional specification derived from a first exploratory pass, which is
identified in \Cref{sec:conclusion} as future work. The requirement as written
is satisfied; the requirement as intended is not, and recording the difference
is more useful than restating the first.

\paragraph{D6. Duplicate findings survive deduplication.} The DataGhurhi
campaign reported three separate major findings for a single control that
failed to open its view, differing only in wording. Deduplication operates over
finding titles and a similarity threshold, and three descriptions of one fault
generated by different runs fall outside it. \emph{Not revised within the
project.} The effect on a stakeholder is over-counting rather than a false
report, so it was ranked below the attribution work; the fix is to deduplicate
against the target element and screen state rather than against the text.

\subsection{The deviation that was not a defect}

One further result deserves its own statement because it looks like a failure
and is not. On broken-workshop, a web application with known faults, the agent
found none. The five completed runs exercised form validation, boundary
lengths, special characters and whitespace handling, and the application
behaved correctly in all five. Whether the faults were outside the reachable
surface, or present and not provoked, cannot be determined from the campaign
record. No revision was made because none is indicated by the evidence: the
run is a demonstration of the measurement gap rather than of a defect in the
system, and it is the single clearest argument in this report for the
seeded-defect experiment that \Cref{sec:conclusion} identifies as the necessary
next piece of work.
